\documentclass[11pt,letterpaper]{article}
\usepackage[margin=0.85in]{geometry}
\usepackage[T1]{fontenc}
\usepackage{amsmath,amssymb,booktabs}
\usepackage{xcolor}
\usepackage{hyperref}
\hypersetup{colorlinks=true,linkcolor=blue!45!black,urlcolor=blue!45!black,
  pdftitle={Camera-to-Printer Frame Transformation}}
\setlength{\parindent}{0pt}
\setlength{\parskip}{4pt}
\setlength{\emergencystretch}{2em}
\newcommand{\p}{\mathbf p}
\newcommand{\R}{\mathbf R}
\newcommand{\T}{\mathbf T}
\newcommand{\tr}{\mathbf t}
\newcommand{\K}{\mathbf K}
\newcommand{\tp}{\mathsf T}

\title{Camera-to-Printer Frame Transformation}
\date{}

\begin{document}
\maketitle
\vspace{-1.5em}

\textbf{The idea.} Print four dots at known printer coordinates. Detect them
in the camera image and calculate the camera pose. Reverse that transformation
to convert predicted 3D surface points into printer coordinates. Then calculate
the nozzle height along the desired path.

\section{Overview}
There are two operations:
\[
\underbrace{\text{Marker}\;\xrightarrow{\;\R,\tr\;}\;\text{Camera}}_
 {\text{3D rotation and translation}}
\quad\xrightarrow{\;\K\;\text{and depth division}\;}\quad
\underbrace{\text{Image}}_{\text{2D pixels}}.
\]
Since we know each dot's printer XYZ, we can express those points directly in the
printer frame, called $P$. The depth-camera frame is called $C$.

The figure's focal lengths $F_x,F_y$ and image offsets $\delta_x,\delta_y$
are written here as $f_x,f_y,c_x,c_y$. They form the intrinsic matrix
\begin{equation}
 \K=\begin{bmatrix}f_x&0&c_x\\0&f_y&c_y\\0&0&1\end{bmatrix}.
\end{equation}
These parameters and lens distortion must already be calibrated for the
selected camera stream. The square is used to find the camera's rotation
and position relative to the printer.

\section{Print and identify the four reference points}
For a square of side $L$ at surface height $Z_0$, the printer coordinates are
\begin{equation}
\begin{aligned}
 \p^P_1&=(X_0,\;Y_0,\;Z_0)^\tp,\\
 \p^P_2&=(X_0+L,\;Y_0,\;Z_0)^\tp,\\
 \p^P_3&=(X_0+L,\;Y_0+L,\;Z_0)^\tp,\\
 \p^P_4&=(X_0,\;Y_0+L,\;Z_0)^\tp.
\end{aligned}
\label{eq:square}
\end{equation}
Detect the corresponding image locations $(u_i,v_i)$, using the same dot
labels. Keep the surface stationary during calibration. $Z_0$ must describe
the observed feature's actual height; it is not automatically the nozzle Z
used to print the dot. All XYZ values and translations below use millimetres.

\section{Calculate printer-to-camera rotation and translation}
Write the rotation and translation explicitly as
\begin{equation}
 \R=\begin{bmatrix}
 r_{11}&r_{12}&r_{13}\\
 r_{21}&r_{22}&r_{23}\\
 r_{31}&r_{32}&r_{33}
 \end{bmatrix},
 \qquad
 \tr=\begin{bmatrix}t_x\\t_y\\t_z\end{bmatrix}.
\end{equation}
The forward transformation is rotation followed by translation:
\begin{equation}
 \underbrace{\begin{bmatrix}x_C\\y_C\\z_C\end{bmatrix}}_{\p_C}
 =
 \underbrace{\begin{bmatrix}
 r_{11}&r_{12}&r_{13}\\
 r_{21}&r_{22}&r_{23}\\
 r_{31}&r_{32}&r_{33}
 \end{bmatrix}}_{\R}
 \underbrace{\begin{bmatrix}X\\Y\\Z\end{bmatrix}}_{\p_P}
 +\underbrace{\begin{bmatrix}t_x\\t_y\\t_z\end{bmatrix}}_{\tr}.
 \label{eq:forward}
\end{equation}
Multiplying each row gives the three coordinate equations:
\begin{equation}
 \begin{aligned}
 x_C&=r_{11}X+r_{12}Y+r_{13}Z+t_x,\\
 y_C&=r_{21}X+r_{22}Y+r_{23}Z+t_y,\\
 z_C&=r_{31}X+r_{32}Y+r_{33}Z+t_z.
 \end{aligned}
 \label{eq:forward-components}
\end{equation}
The camera projects this point into the image. For distortion-corrected
pixels, the complete equation $z_C(u,v,1)^\tp=\K(\R\p_P+\tr)$ expands to
\begin{equation}
 \boxed{z_C\begin{bmatrix}u\\v\\1\end{bmatrix}
 =
 \begin{bmatrix}f_x&0&c_x\\0&f_y&c_y\\0&0&1\end{bmatrix}
 \begin{bmatrix}
 r_{11}&r_{12}&r_{13}&t_x\\
 r_{21}&r_{22}&r_{23}&t_y\\
 r_{31}&r_{32}&r_{33}&t_z
 \end{bmatrix}
 \begin{bmatrix}X\\Y\\Z\\1\end{bmatrix}.}
 \label{eq:projection}
\end{equation}
Dividing by $z_C$ gives the pixel coordinates directly:
\begin{equation}
 \boxed{\begin{aligned}
 u&=f_x\frac{r_{11}X+r_{12}Y+r_{13}Z+t_x}
                 {r_{31}X+r_{32}Y+r_{33}Z+t_z}+c_x,\\[6pt]
 v&=f_y\frac{r_{21}X+r_{22}Y+r_{23}Z+t_y}
                 {r_{31}X+r_{32}Y+r_{33}Z+t_z}+c_y.
 \end{aligned}}
 \label{eq:pixel-components}
\end{equation}
For each reference point, substitute its known $X_i,Y_i,Z_i$ and compare
the calculated pixels with the detected $u_i,v_i$. All four points share
the same rotation and translation.
The factor $z_C$ is essential: a 3D point must be divided by its depth to
obtain a pixel location. A pixel alone does not determine a 3D point.

Pass the four known printer points, their image locations, camera intrinsics,
and distortion parameters to a PnP solver. OpenCV's \texttt{solvePnP} returns
the rotation vector and translation from the supplied object frame to the
camera frame~\cite{pnp}. Here that object frame is $P$:
\begin{equation}
 (\texttt{rvec},\texttt{tvec})
 =\operatorname{solvePnP}\!\left(
 \{\p_i^P\},\{(u_i,v_i)\},\K,\text{distortion}\right).
\end{equation}
Convert \texttt{rvec} into $\R$ using Rodrigues' formula; $\tr=\texttt{tvec}$.
Use a planar solver appropriate for the point order. Check its pose against
measured depth and the known camera orientation, because planar image data
can admit competing poses. If dots are detected in a separate RGB camera,
also use its calibrated relationship to the depth camera.

\section{Reverse the transformation}
For printing we need camera-to-printer coordinates. Rearranging
Equation~\eqref{eq:forward} gives
\begin{equation}
 \boxed{\p_P=\R^\tp(\p_C-\tr).}
 \label{eq:inverse}
\end{equation}
The transpose is the inverse of a rotation matrix. The camera-to-printer
matrix, written entry by entry, is
\begin{equation}
 \boxed{\T_{P\leftarrow C}=
 \begin{bmatrix}
 r_{11}&r_{21}&r_{31}&-(r_{11}t_x+r_{21}t_y+r_{31}t_z)\\
 r_{12}&r_{22}&r_{32}&-(r_{12}t_x+r_{22}t_y+r_{32}t_z)\\
 r_{13}&r_{23}&r_{33}&-(r_{13}t_x+r_{23}t_y+r_{33}t_z)\\
 0&0&0&1
 \end{bmatrix}.}
 \label{eq:matrix}
\end{equation}
It acts on a camera point as
\begin{equation}
 \begin{bmatrix}X\\Y\\Z\\1\end{bmatrix}
 =\T_{P\leftarrow C}\begin{bmatrix}x_C\\y_C\\z_C\\1\end{bmatrix}.
\end{equation}
Equivalently, the individual printer coordinates are
\begin{equation}
 \boxed{\begin{aligned}
 X&=r_{11}(x_C-t_x)+r_{21}(y_C-t_y)+r_{31}(z_C-t_z),\\
 Y&=r_{12}(x_C-t_x)+r_{22}(y_C-t_y)+r_{32}(z_C-t_z),\\
 Z&=r_{13}(x_C-t_x)+r_{23}(y_C-t_y)+r_{33}(z_C-t_z).
 \end{aligned}}
 \label{eq:inverse-components}
\end{equation}

\newpage
\section{Transform the predicted surface}
The existing height map uses a metric camera XY grid. A valid predicted cell
at row $i$, column $j$, and future time $t_a$ represents the 3D point
\begin{equation}
 \widehat\p_C(i,j;t_a)=
 \begin{bmatrix}x_j\\y_i\\\widehat h_C(i,j;t_a)\end{bmatrix},
 \quad
 \begin{aligned}
 x_j&=x_{\min}+(j+\tfrac12)\Delta,\\
 y_i&=y_{\min}+(i+\tfrac12)\Delta,
 \end{aligned}
\end{equation}
where $\Delta$ is the cell size and indices start at zero. Transform this
complete point using the saved calibration:
\begin{equation}
 \boxed{\widehat\p_P(i,j;t_a)=
 \R^\tp\!\left(
 \begin{bmatrix}x_j\\y_i\\\widehat h_C(i,j;t_a)\end{bmatrix}-\tr\right).}
 \label{eq:prediction}
\end{equation}
Here $x_j,y_i$ are already in millimetres, so no pixel deprojection is needed.
Use absolute, unnormalized camera heights. If the model predicts a height
change, add its reference height first.

Transform X, Y, and Z together: with a tilted camera, changing camera depth
can also change printer X and Y. The transformed points generally form an
irregular XY grid. Interpolate or rasterize them onto a printer XY grid to
obtain $\widehat h_P(X,Y;t_a)$, retaining invalid regions. This assumes the
surface has one relevant height at each printer XY location.

\section{Calculate the nozzle height}
Choose the desired printer path points $(X_k,Y_k)$. Predict the surface for
the time $t_k$ when the nozzle will reach each point. For vertical gap $g$,
\begin{equation}
 \boxed{Z_k=\widehat h_P(X_k,Y_k;t_k)+g,\qquad
 \p_{\mathrm{nozzle},k}^P=\begin{bmatrix}X_k\\Y_k\\Z_k\end{bmatrix}.}
 \label{eq:nozzle}
\end{equation}

\begin{thebibliography}{1}
\bibitem{pnp} OpenCV, \emph{Perspective-n-Point (PnP) pose computation}.
\url{https://docs.opencv.org/4.x/d5/d1f/calib3d_solvePnP.html}
\end{thebibliography}
\end{document}
